\problemname{Poplåtar}
Your band \emph{The Length-First Searchers} has written a new song.
The music is enchanting and the lyrics are catchy, so it might very well top the charts.

The only thing remaining before you and your band can go to the studio and record the song is to combine the music with the lyrics.
You have already decided that the song should start with an intro, followed by a chorus, and end with an outro (nobody likes verses these days anyway).
During each part of the song, you will sing a portion of the lyrics.
During the intro and outro you can sing whatever you want, but during the chorus you must sing in a more structured way: one line of text at a time.
To match the rhythm of the music, each line of the chorus must also have the same length.
You do not actually have to sing anything during the intro and outro, but the chorus must consist of at least two lines.

There can be very many ways to choose which parts of the lyrics belong to each section, so instead of choosing one at random, you decide to maximize the \emph{rhyme value} of the chorus.
The rhyme value is defined as follows: let $N$ be the number of lines in the chorus and $S$ be the length of the longest common suffix that all lines share.
Then the rhyme value equals $N \cdot S$.

\section*{Input}
The first and only line of input consists of a string, the song's lyrics.
The string contains only lowercase letters (\texttt{a-z}) and underscores (``\texttt{\_}'').
Its length is between $2$ and $15\,000$.

\section*{Output}
Print a single integer: the largest possible rhyme value you can achieve.

\section*{Scoring}
Your solution will be tested on a set of test groups.
To earn points for a group, you must pass all test cases in that group.

\noindent
\begin{tabular}{| l | l | l |}
\hline
Group & Points & Constraints \\ \hline
$1$   & $30$       & The string's length is at most $400$ \\ \hline
$2$   & $30$       & The string's length is at most $4000$ \\ \hline
$3$   & $40$       & The string's length is at most $15\,000$ \\ \hline
\end{tabular}
